\chapter{Integration}

Integration is the inverse operation of differentiation. While differentiation measures the instantaneous slope of a curve, integration calculates the continuous accumulation of quantities, geometrically represented as the area under a curve. In data science, integration is the mathematical cornerstone for probability. Specifically, it is used to calculate cumulative probability distributions, find expected values (means) of continuous random variables, and determine expected loss over distributions.

\section{Riemann Sums and the Definite Integral}

Consider a continuous function $f(x)$ defined on a closed interval $[a, b]$. We want to find the exact area bounded by the curve $y = f(x)$, the x-axis, and the vertical boundaries $x = a$ and $x = b$.

\begin{definitionbox}{Riemann Sums and Definite Integral}
\begin{itemize}
    \item \textbf{Partition:} Divide the interval $[a, b]$ into $n$ subintervals of equal width:
    $$\Delta x = \frac{b - a}{n}$$
    \item \textbf{Riemann Sum:} Approximate the area using $n$ rectangles. The sum of the areas of these rectangles is:
    $$\text{Riemann Sum} = \sum_{i=1}^n f(x_i^*) \Delta x$$
    where $x_i^*$ is a sample point in the $i$-th subinterval $[x_{i-1}, x_i]$.
    \item \textbf{Definite Integral:} The definite integral of $f$ from $a$ to $b$ is defined as the limit of the Riemann sum as the number of rectangles approaches infinity ($n \to \infty$):
    $$\int_{a}^{b} f(x) \dx = \limit{n}{\infty} \sum_{i=1}^n f(x_i^*) \Delta x$$
    The numbers $a$ and $b$ are called the lower and upper limits of integration.
\end{itemize}
\end{definitionbox}

\textbf{Data Science Relevance:} In statistics, a Probability Density Function (PDF), $p(x)$, describes the likelihood of a continuous random variable taking on a certain value. The probability that the variable falls between $a$ and $b$ is exactly the area under the curve between those points, calculated as $\int_a^b p(x) \dx$. 



\section{The Fundamental Theorem of Calculus (FTC)}

The FTC is one of the most remarkable discoveries in mathematics because it connects the two seemingly unrelated concepts of differentiation (slopes) and integration (areas).

\begin{theorembox}{The Fundamental Theorem of Calculus}
Let $f(x)$ be a continuous function on $[a, b]$.
\begin{itemize}
    \item \textbf{FTC Part 1 (Derivative of an Integral):} If we define an area accumulation function $g(x) = \int_{a}^{x} f(t) \,dt$, then its rate of change is simply the curve's height:
    $$g'(x) = \deriv{}{x} \left[ \int_{a}^{x} f(t) \,dt \right] = f(x)$$
    \item \textbf{FTC Part 2 (Evaluation Theorem):} If $F(x)$ is any antiderivative of $f(x)$ (meaning $F'(x) = f(x)$), then the definite integral can be evaluated simply by subtracting the bounds:
    $$\int_{a}^{b} f(x) \dx = F(b) - F(a) = \Big[ F(x) \Big]_a^b$$
\end{itemize}
\end{theorembox}



\section{Indefinite Integrals and Basic Techniques}

An indefinite integral represents the family of all possible antiderivatives of a function. Because the derivative of a constant is zero, integration leaves ambiguity that we capture with $+ C$.
$$\int f(x) \dx = F(x) + C$$

\begin{formulabox}{Standard Integration Formulas}
\begin{itemize}
    \item \textbf{Power Rule:} $\int x^n \dx = \frac{x^{n+1}}{n+1} + C \quad (\text{for } n \neq -1)$
    \item \textbf{Log Rule:} $\int \frac{1}{x} \dx = \ln|x| + C$
    \item \textbf{Exponential Rule:} $\int e^x \dx = e^x + C$
    \item \textbf{Constant Multiple Rule:} $\int c \cdot f(x) \dx = c \int f(x) \dx$
    \item \textbf{Sum Rule:} $\int [f(x) \pm g(x)] \dx = \int f(x) \dx \pm \int g(x) \dx$
\end{itemize}
\end{formulabox}

\begin{examplebox}{Definite Integral Evaluation}
Evaluate the definite integral: $\int_{1}^{3} (3x^2 + 2x - 1) \dx$.
\begin{enumerate}
    \item Find the general antiderivative $F(x)$ using the Sum and Power Rules:
    $$F(x) = \int (3x^2 + 2x - 1) \dx = 3\left(\frac{x^3}{3}\right) + 2\left(\frac{x^2}{2}\right) - x = x^3 + x^2 - x$$
    \item Evaluate $F(x)$ at the upper and lower limits:
    $$F(3) = (3)^3 + (3)^2 - 3 = 27 + 9 - 3 = 33$$
    $$F(1) = (1)^3 + (1)^2 - 1 = 1 + 1 - 1 = 1$$
    \item Apply FTC Part 2 ($F(b) - F(a)$):
    $$\int_{1}^{3} (3x^2 + 2x - 1) \dx = F(3) - F(1) = 33 - 1 = 32$$
\end{enumerate}
\end{examplebox}

\subsection{The Substitution Rule (U-Substitution)}
U-substitution is the "reverse Chain Rule". If a function and its derivative are both present inside the integral, we can compress it.
\begin{theorembox}{The Substitution Rule}
If $u = g(x)$ is a differentiable function, then $du = g'(x) \dx$, allowing the transformation:
$$\int f(g(x))g'(x) \dx = \int f(u) \,du$$
\end{theorembox}

\begin{examplebox}{Substitution Rule Application}
Evaluate the indefinite integral: $\int 2x e^{x^2} \dx$.
\begin{enumerate}
    \item Let the inner exponent $u = x^2$. Differentiating with respect to $x$ gives:
    $$du = 2x \dx$$
    \item Substitute $u$ and $du$ entirely into the original integral:
    $$\int e^{x^2} \cdot (2x \dx) = \int e^u \,du$$
    \item Integrate with respect to $u$:
    $$\int e^u \,du = e^u + C$$
    \item Substitute back $u = x^2$ to return to the original variable:
    $$e^{x^2} + C$$
\end{enumerate}
\end{examplebox}

\section{Area Between Curves}

If $f(x) \ge g(x)$ for all $x$ in $[a, b]$, the area $A$ of the region bounded by the curves $y = f(x)$ on top, $y = g(x)$ on the bottom, and the vertical lines $x = a$, $x = b$ is simply the integral of the difference:
$$A = \int_{a}^{b} [f(x) - g(x)] \dx$$

\newpage
\section{Practice Exercises}
\textit{Try to solve these exercises independently. Detailed step-by-step solutions are available in the Appendix at the end of the book.}

\begin{enumerate}
    \item \textbf{Pure Math:} Evaluate the indefinite integral: $\int \left( 4x^3 - \frac{1}{x} + e^x \right) \dx$.
    \item \textbf{Pure Math:} Use u-substitution to evaluate the integral: $\int x^2(x^3 + 5)^4 \dx$.
    \item \textbf{Data Science:} The Probability Density Function of a uniform distribution is given by $p(x) = 0.5$ on the interval $[0, 2]$. Verify that the total probability over this entire interval is exactly $1$ by calculating $\int_{0}^{2} 0.5 \dx$.
    \item \textbf{Pure Math:} Evaluate the definite integral using FTC Part 2: $\int_{1}^{2} \frac{1}{x^2} \dx$.
    \item \textbf{Pure Math:} Find the area of the region enclosed between the curves $y = x^2$ and $y = x$ on the interval $[0, 1]$.
\end{enumerate}